% ============================================================
%  SECCIÓN 3 – PERFILES DE USUARIO
% ============================================================
\section{Perfiles de Usuario}

% ──────────────────────────────────────────────────────────
\subsection{Perfil 1: Donante Potencial}

\subsubsection{Datos Demográficos}
Es una persona entre 18 y 65 años que vive en Popayán o sus alrededores. Cuenta con documento de identidad y suele movilizarse por estudio, trabajo o diligencias dentro de la ciudad. Su peso es mayor a 50\,kg.

\subsubsection{Objetivos y Necesidades}
Quiere ayudar donando sangre, pero necesita que le expliquen de forma sencilla si cumple los requisitos y qué debe hacer antes de ir. También necesita saber con claridad dónde puede donar, si hay una jornada cerca, y poder agendar una hora para no perder tiempo.

\subsubsection{Motivaciones}
Le motiva salvar vidas, apoyar a su comunidad y sentir que su acción tiene un impacto real. Responde mejor cuando entiende que lo necesitan y que el proceso es rápido y seguro.

\subsubsection{Frustraciones o Problemas}
Se frustra cuando la información está dispersa en páginas web o redes sociales y no sabe cuál es la fuente oficial o actualizada. Se desanima si llega a donar y le informan en el punto que no puede hacerlo por algo que nadie le explicó antes. También le generan dudas los mitos, por ejemplo creer que donar lo debilita, afecta su peso o su vida sexual, o que la sangre se vende.

\subsubsection{Habilidades Tecnológicas}
Usa el celular a diario para redes sociales, mensajería y trámites básicos. Puede registrarse en una aplicación si el proceso es corto y el lenguaje es claro.

\subsubsection{Contexto de Uso}
Normalmente revisa el teléfono en ratos libres. Si recibe una alerta de necesidad o una jornada de donación, decide rápido si puede ir, pero necesita saber el tiempo estimado y tener confirmación.

\subsubsection{Comportamientos y Hábitos Relevantes}
Si es hombre puede donar cada tres meses y si es mujer cada seis meses, por lo que necesita que la aplicación le recuerde cuándo vuelve a ser apto. Puede tener dudas sobre temas sensibles del requisito (por ejemplo, cambios de pareja sexual o antecedentes de hepatitis) y por eso requiere un cuestionario privado, respetuoso y fácil de entender.

% ──────────────────────────────────────────────────────────
\subsection{Perfil 2: Solicitante de Sangre}

\subsubsection{Datos Demográficos}
Persona adulta (familiar, amigo o cuidador de un paciente) que reside en Popayán o se encuentra temporalmente en la ciudad por motivos de atención médica. Puede pertenecer a cualquier estrato socioeconómico y nivel de formación académica.

\subsubsection{Objetivos y Necesidades}
Necesita conseguir donantes de un tipo de sangre específico en el menor tiempo posible para un paciente que se encuentra en un centro de salud. Requiere una herramienta que le permita publicar la necesidad, llegar a personas compatibles cercanas y coordinar la logística de la donación de manera clara y segura.

\subsubsection{Motivaciones}
La urgencia de la situación médica de un ser querido lo impulsa a actuar con rapidez. Valora cualquier medio que le permita ampliar su alcance más allá de su círculo social inmediato.

\subsubsection{Frustraciones o Problemas}
Las cadenas de mensajes en redes sociales o aplicaciones de mensajería no le permiten filtrar por compatibilidad ni por cercanía. No tiene forma de saber quién realmente puede o quiere donar. La incertidumbre y la espera generan angustia adicional en un contexto ya estresante.

\subsubsection{Habilidades Tecnológicas}
Similares al donante potencial; usa el celular para comunicaciones básicas. En situación de urgencia, necesita que el proceso de publicación sea extremadamente rápido e intuitivo.

\subsubsection{Contexto de Uso}
Se encuentra en un centro de salud o cerca de él, con acceso a su teléfono móvil. Necesita publicar la solicitud, recibir respuestas y coordinar con donantes, todo desde la misma herramienta.

\subsubsection{Comportamientos y Hábitos Relevantes}
En situación de emergencia, su capacidad de atención y paciencia puede verse reducida. Por ello, la aplicación debe minimizar los pasos requeridos y ofrecer retroalimentación constante sobre el estado de su solicitud (donantes notificados, donantes aceptados, etc.).

% ──────────────────────────────────────────────────────────
\subsection{Perfil 3: Profesional de Salud}

\subsubsection{Datos Demográficos}
Profesional vinculado a una institución de salud (hospital, clínica, banco de sangre, universidad con facultad de medicina) en Popayán o sus alrededores. Puede ser médico, enfermero, coordinador de banco de sangre o personal administrativo autorizado.

\subsubsection{Objetivos y Necesidades}
Busca una herramienta que le permita crear y publicar campañas de donación para su institución, alcanzando a la mayor cantidad posible de potenciales donantes en la zona. Necesita gestionar cupos, comunicar requisitos y fechas, y hacer seguimiento de las reservas realizadas.

\subsubsection{Motivaciones}
Contribuir al abastecimiento de sangre en su institución y en la comunidad. Le motiva contar con un canal digital directo que complemente los métodos tradicionales de convocatoria.

\subsubsection{Frustraciones o Problemas}
Los medios tradicionales de convocatoria (volantes, radio, redes sociales institucionales) tienen alcance limitado y no permiten medir respuestas ni organizar la logística de donantes. Las campañas exitosas dependen de factores no controlables, como la viralización espontánea de la información.

\subsubsection{Habilidades Tecnológicas}
Posee habilidades digitales medias a altas. Está acostumbrado a utilizar sistemas de información hospitalaria y aplicaciones profesionales.

\subsubsection{Contexto de Uso}
Planifica las campañas desde su lugar de trabajo, con tiempo suficiente para completar formularios detallados. Consulta el estado de las reservas y las estadísticas de la campaña de forma periódica.

\subsubsection{Comportamientos y Hábitos Relevantes}
Valora la formalidad y el respaldo institucional. Necesita que la plataforma le permita publicar bajo el nombre de su institución y que la campaña se muestre como contenido verificado y confiable para la comunidad.
